\section{A finite tag clock}
\label{sec:clock}

A prefix and a suffix may choose different metadata. We need a finite
residue test that either makes their choices agree or confines every
successful disagreement to one fixed input--output pair. A deterministic
affine correction can absorb the latter case without changing the linear
part to be recovered. We also need many possible cuts: the number spoiled
by the clock must remain bounded when more copies are added.

The following lemma supplies both properties. Its tags are finite metadata,
never input letters. Its roster is the finite index set of possible cut
copies. The modulus may grow with this roster, but the obstruction bound
does not. Throughout this section, $\log$ denotes the natural logarithm.

\begin{lemma}[Tag clock]
\label{lem:clock}
Let $\Lambda$ be a nonempty finite set, let $m=|\Lambda|$, let
$\mathcal R$ be any finite set, and let $B>0$.  There are a positive integer
$n$, sets $I_\alpha,J_\alpha\subseteq\mathbb Z/n\mathbb Z$ for
$\alpha\in\Lambda$, a set $W\subseteq\mathbb Z/n\mathbb Z$, and residues
$v_p\in\mathbb Z/n\mathbb Z$ for $p\in\mathcal R$, with the following
properties.
\begin{enumerate}[label=\textup{(\roman*)}]
\item Every prime factor of $n$ exceeds $B$.
\item For every $\ell\in\mathbb Z/n\mathbb Z$, the directed graph
\[
 \mathcal E_\ell
 =\{(\alpha,\beta)\in\Lambda^2:
                 \ell\in I_\alpha+J_\beta\}
\]
is either contained in the diagonal
$\{(\alpha,\alpha):\alpha\in\Lambda\}$ or consists of exactly one
ordered pair $(\alpha,\beta)$ with $\alpha\ne\beta$.
\item If $\ell\notin W$, then $\mathcal E_\ell$ is the full diagonal.
\item For every $u\in\mathbb Z/n\mathbb Z$,
\begin{equation}
 \left|\left\{(p,p')\in\mathcal R^2:
       p\ne p',\quad u+v_{p'}-v_p\in W\right\}\right|
 \le c_{\mathrm{cl}}\log(2m),
 \label{eq:clock-roster}
\end{equation}
where the absolute constant $c_{\mathrm{cl}}=1\,300\,000$ is valid.
\end{enumerate}
In particular, $c_{\mathrm{cl}}$ is independent of $B$, $m$, and
$|\mathcal R|$.
\end{lemma}

The ordered-pair estimate will control false guesses of a word's origin.
Choose integer cut offsets congruent to their assigned roster residues.
If an actual cut has offset $p$ from an origin $t$, interpreting it as
offset $p'$ instead gives the candidate origin $t+p-p'$. For a fixed end
anchor $q$, its total residue is $q-t+p'-p$. Property~\textup{(iv)}
therefore bounds how many such wrong candidates have a bad total residue,
even when the roster contains many copies of every cut kind. The modulus
may be very large; the number of bad pairs is the quantity needed below.

\begin{proof}
We first specify finitely many cyclic coordinates and construct the sets
and roster values in their product.  At the end we identify this product
with a single cyclic group.  Put
\[
 a=\lceil64\log(2m)\rceil,\qquad
 b=\lceil\log_2 m\rceil,\qquad K_{\mathrm{cl}}=400,
 \qquad r=|\mathcal R|.
\]
Here $K_{\mathrm{cl}}$ is fixed once and for all.

\paragraph{Partitions and coordinate layout.}
Choose $a$ assignments of $\Lambda$ to a left side and a right side,
allowing either side to be empty, so that, whenever
\begin{equation}
 \alpha\ne\beta,\qquad
       \{\alpha,\beta\}\ne\{\gamma,\delta\},
 \label{eq:clock-separation-constraint}
\end{equation}
some partition puts $\alpha$ on the left, $\beta$ on the right, and
$\gamma,\delta$ on the same side.  We allow $\gamma=\delta$.
To prove existence, independently assign each tag to a side with
probability $1/2$.  The required assignment is feasible: if
$\{\gamma,\delta\}$ contains $\alpha$, put both on the left; if it
contains $\beta$, put both on the right; and if it contains neither,
put both on the left.  The excluded equality in
Equation~\eqref{eq:clock-separation-constraint} ensures that these prescriptions
never force $\alpha$ and $\beta$ onto the same side.  Specifying sides
for at most four distinct tags has probability at least $1/16$.
There are at most $m^4$ constraints.  With $a$ independent partitions,
the expected number that all partitions miss is at most
\[
 m^4(15/16)^a\le m^4\exp(-a/16)
                  \le m^4(2m)^{-4}=1/16<1.
\]
The number missed is a nonnegative integer, so some choice misses none.
We give each of these $a$ partitions one coordinate, called a
\emph{first coordinate}.

Also choose $b$ partitions into Type~1 and Type~2 that separate every
two distinct tags.  For example, label the tags by distinct binary
strings of length $b$, and use the bit at each position.  When $m=1$,
$b=0$ and there are no such partitions.
Give each Type partition a batch of coordinates as follows.  List every
ordered list of $K_{\mathrm{cl}}$ directed pairs $(p,p')$ of distinct
roster labels whose underlying unordered edges are distinct and form a
forest.  A forest here is a finite simple graph with no cycles.  Give
each such list one dedicated coordinate in the batch, and add spare
coordinates if needed to make the batch size positive and odd.  All
these lists form a finite set.  This defines a finite coordinate set
$\mathcal C$ before choosing any prime orders.

Choose pairwise distinct primes $N_c$, $c\in\mathcal C$, satisfying
\begin{equation}
 N_c>\max\{B,400,2^{r+1}\}.
 \label{eq:clock-prime-threshold}
\end{equation}
This is possible recursively: the product of all primes at most any
given finite bound, plus one, has a prime divisor exceeding that bound.
At each step include the previously chosen primes in the bound.
Write
\[
 G=\prod_{c\in\mathcal C}\mathbb Z/N_c\mathbb Z.
\]
For $x\in\mathbb Z/N_c\mathbb Z$, identify $x/N_c$ with its point on
$\mathbb T=\mathbb R/\mathbb Z$.  On this circle write
$\|z\|_{\mathbb T}=\min_{j\in\mathbb Z}|z-j|$ and
$\operatorname{dist}_{\mathbb T}(z,A)
=\min_{a'\in A}\|z-a'\|_{\mathbb T}$ for a nonempty finite set $A$.

\paragraph{The first coordinates.}
In the coordinate for a left/right partition, impose the requirements
\[
\begin{array}{c|cc}
 \text{side of }\alpha & I_\alpha\text{ coordinate}
                         &J_\alpha\text{ coordinate}\\ \hline
 \text{left} &\{0\}&(\mathbb Z/N_c\mathbb Z)\setminus\{0\}\\
 \text{right}&(\mathbb Z/N_c\mathbb Z)\setminus\{0\}&\{0\}.
\end{array}
\]
Any nonzero total is a sum for a diagonal pair: put the total in the
required nonzero summand and put zero in the other.  At total zero,
however, no edge between two tags on the same side is possible, since
it would have exactly one nonzero summand.

Suppose an off-diagonal edge $\alpha\to\beta$ occurs at a total
$\ell\in G$.  For any $(\gamma,\delta)$ satisfying
Equation~\eqref{eq:clock-separation-constraint}, choose its separating partition.
The edge $\alpha\to\beta$ forces total zero in that coordinate, because
its two summands are both required to be zero.  The competing edge
$\gamma\to\delta$ is impossible there because its tags are on the
same side.  Thus the only other edge that could coexist with
$\alpha\to\beta$ is $\beta\to\alpha$.  This argument also excludes
every diagonal edge, because the constraints explicitly include
$\gamma=\delta$.

\paragraph{Type~1 batches: all sums and separated differences.}
Put
\[
 A=\{0,1,3\}/6\subseteq\mathbb T,
 \qquad D=\{2,4,5\}/6\subseteq\mathbb T.
\]
For a tag of Type~1 in a batch, require in every coordinate $c$ of that
batch
\[
 \operatorname{dist}_{\mathbb T}(i_c/N_c,A)\le1/16
 \quad\text{for }i\in I_\alpha,
 \qquad
 \operatorname{dist}_{\mathbb T}(j_c/N_c,D)\le1/16
 \quad\text{for }j\in J_\alpha.
\]
The sums of the two center sets contain every sixth of the circle:
\[
 \{0,1,3\}+\{2,4,5\}=\mathbb Z/6\mathbb Z.
\]
Every target $z\in\mathbb T$ is therefore within $1/12$ of some center
sum $a'+d'$.  Choose a signed adjustment $\theta$, with
$|\theta|\le1/12$, such that $z=a'+d'+\theta$ on the circle.  The
continuous summands $a'+\theta/2$ and $d'+\theta/2$ are each within
$1/24$ of their respective centers.  If $z$ is a point of the
$N_c$-grid, round the first summand to the nearest grid point and define
the second to be $z$ minus the rounded first.  Both summands then lie
on the grid and each changes by circular distance at most $1/(2N_c)$.
Since $N_c>24$,
\[
 \frac1{24}+\frac1{2N_c}<\frac1{16}.
\]
Thus the Type~1 sumset is the entire batch group, with the choices made
independently in its coordinates.

The two center sets have circular distance at least $1/6$ from each
other.  Consequently, every Type~1 difference $i-j$ satisfies
\begin{equation}
 \|(i_c-j_c)/N_c\|_{\mathbb T}
 \ge \frac16-\frac2{16}=\frac1{24}
 \quad\text{in every coordinate of its batch}.
 \label{eq:clock-type-one-difference}
\end{equation}

\paragraph{Type~2 batches: majority sums and differences.}
Suppose a batch has $2s+1$ coordinates.  Call a coordinate value
\emph{small} if its circular distance from zero is at most $1/100$.
For a tag of Type~2, both $I_\alpha$ and $J_\alpha$ require a strict
majority, that is, at least $s+1$, of the batch coordinates to be small.
There are no restrictions on the other coordinates.

Their sumset contains every target vector $z$ with at least one small
coordinate.  Indeed, at one such coordinate set the first summand equal
to $z$ and the second to zero, making both small.  Partition the other
$2s$ coordinates into two sets of size $s$.  On the first set make the
first summand zero, and on the second set make the second summand zero;
in each coordinate determine the remaining summand by the target.
Both vectors have at least $s+1$ small coordinates.  Conversely, the
sets of small coordinates of two vectors satisfying the majority
condition must intersect.  Every Type~2 difference therefore satisfies
\begin{equation}
 \|(i_c-j_c)/N_c\|_{\mathbb T}\le2/100
 \quad\text{for at least one coordinate in its batch}.
 \label{eq:clock-type-two-difference}
\end{equation}
This argument also covers a one-coordinate batch, with $s=0$.

\paragraph{The exclusive graph alternative.}
Define each $I_\alpha\subseteq G$ and $J_\alpha\subseteq G$ by imposing
all its first-coordinate and batch requirements.  These requirements
concern disjoint coordinate blocks.  We already know that an
off-diagonal edge can coexist only with its reversal.  If both
$\alpha\to\beta$ and $\beta\to\alpha$ occurred, there would be elements
$i_\alpha\in I_\alpha$, $j_\alpha\in J_\alpha$,
$i_\beta\in I_\beta$, and $j_\beta\in J_\beta$ with
\[
 i_\alpha+j_\beta=\ell=i_\beta+j_\alpha,
 \qquad i_\alpha-j_\alpha=i_\beta-j_\beta.
\]
Choose a Type partition separating $\alpha$ and $\beta$.  In its batch
the common difference would satisfy the bounds in both
Equation~\eqref{eq:clock-type-one-difference} and
Equation~\eqref{eq:clock-type-two-difference}, which is impossible because
$2/100<1/24$.  This proves the alternative in \textup{(ii)} in $G$.

Let $W\subseteq G$ be the union of the following obstructions:
\begin{enumerate}[label=\textup{(\alph*)}]
\item a first coordinate of the total is zero;
\item in some Type batch, none of the total's coordinates is small.
\end{enumerate}
If $\ell\notin W$, every diagonal pair has a decomposition of $\ell$
in every first coordinate and every batch, by the sumset arguments
above.  Combining these decompositions gives
$\ell\in I_\alpha+J_\alpha$ for every $\alpha$.  The exclusive graph
alternative then implies that the graph is the full diagonal.  This
proves \textup{(iii)}.  No converse assertion about $W$ is needed.

We have now established the graph alternative and the availability of all
diagonal pairs outside $W$. It remains to choose the roster residues so
that every translate produces few ordered pairs in $W$. The first
coordinates and the Type batches require different counting arguments.

\paragraph{Roster values in the first coordinates.}
Enumerate $\mathcal R=\{p_0,\ldots,p_{r-1}\}$.  In every first
coordinate assign
\[
 (v_{p_j})_c=2^j\pmod{N_c}.
\]
The nonzero ordered integer differences are all distinct.  For a
positive difference with $j>i$, the factorization
\[
 2^j-2^i=2^i(2^{j-i}-1)
\]
determines $i$ as the exponent of the largest power of two dividing the difference,
and then determines $j$.  The sign distinguishes the reverse ordering.
For $r\ge2$, all signed differences lie between
$-(2^{r-1}-1)$ and $2^{r-1}-1$; by
Equation~\eqref{eq:clock-prime-threshold}, their range has length less than
$N_c$.  Thus they remain distinct modulo $N_c$.
For any fixed $u\in G$, at most one ordered pair $p\ne p'$ can satisfy
\[
 u_c+(v_{p'})_c-(v_p)_c=0.
\]
This also holds when $r\le1$, since there are no such pairs.

\paragraph{Roster values in a dedicated coordinate.}
Consider a coordinate $c$ dedicated to a list with directed pairs
\[
 (\pi_j,\pi'_j),\qquad 0\le j<K_{\mathrm{cl}}.
\]
Choose $d_j\in\mathbb Z/N_c\mathbb Z$ by rounding $j/K_{\mathrm{cl}}$
to the nearest point of the $N_c$-grid, so that
\[
 \|d_j/N_c-j/K_{\mathrm{cl}}\|_{\mathbb T}
 \le1/(2N_c).
\]
Prescribe $(v_{\pi'_j})_c-(v_{\pi_j})_c=d_j$.  To realize these prescriptions,
root each tree component and assign its root value zero.  On traversing
an edge to a previously unassigned vertex, add $d_j$ if traversing
from $\pi_j$ to $\pi'_j$, and subtract $d_j$ if traversing in the opposite
direction.  The absence of cycles means that each new vertex is
assigned once, with no consistency condition left to check.  Give
isolated roster labels value zero.  Spare coordinates can have arbitrary
roster values, for example all zero.  Distinct roster labels may receive
the same value in a batch coordinate; only the prescribed edge
differences will be used.

For every $u_c\in\mathbb Z/N_c\mathbb Z$, one of the equally spaced
points $j/K_{\mathrm{cl}}$ is within $1/(2K_{\mathrm{cl}})$ of
$-u_c/N_c$.  For that edge,
\begin{equation}
 \bigl\|(u_c+(v_{\pi'_j})_c-(v_{\pi_j})_c)/N_c\bigr\|_{\mathbb T}
 \le\frac1{2K_{\mathrm{cl}}}+\frac1{2N_c}
 <\frac1{100},
 \label{eq:clock-forest-net}
\end{equation}
using $K_{\mathrm{cl}}=400$ and $N_c>400$.  Thus every translate makes
at least one edge of this dedicated forest small in this coordinate.

\paragraph{The bound for one batch.}
Fix $u\in G$ and a batch, and let $S$ be the set of ordered pairs
$p\ne p'$ for which $u+v_{p'}-v_p$ has no small coordinate in that
batch.  Suppose that $|S|>4K_{\mathrm{cl}}^2$.
Forget directions and identify the two orientations of each unordered
edge to obtain a simple graph on the nonisolated roster labels.
This graph contains a forest with
$K_{\mathrm{cl}}$ edges.  To see this, take a spanning forest, obtained
by growing a tree in each connected component until all its vertices
are reached.  If its edge count $f$ were less than $K_{\mathrm{cl}}$,
then, writing $v$ for its number of vertices and $t$ for its number of
components, we would have
\[
 f=v-t,\qquad t\le v/2,\qquad v\le2f<2K_{\mathrm{cl}}.
\]
Here each component has at least two vertices because isolated
vertices were removed.  But then
$|S|\le v(v-1)<4K_{\mathrm{cl}}^2$, a contradiction.

Choose $K_{\mathrm{cl}}$ forest edges, choose for each an orientation
belonging to $S$, and order the resulting directed pairs.  The batch
already has a coordinate dedicated to this list.  By
Equation~\eqref{eq:clock-forest-net}, one of its pairs has a small translated
difference in that coordinate, contradicting membership in $S$.
Consequently every batch contributes at most $4K_{\mathrm{cl}}^2$
ordered pairs to the obstruction.  The forest may depend on $u$;
every possible directed list was allocated a coordinate in advance.
If no such forest can be formed from the roster, the spanning-forest
count alone gives the required bound, and the spare coordinates do
not affect it.

\paragraph{Combining the estimates and the cyclic coordinates.}
Taking a union over the $a$ first coordinates and $b$ batches gives
\[
 \left|\{(p,p'):p\ne p',\ u+v_{p'}-v_p\in W\}\right|
 \le a+4K_{\mathrm{cl}}^2b.
\]
Since $\log 2>1/2$,
\[
 a\le66\log(2m),\qquad
 b\le\frac{\log(2m)}{\log2}\le2\log(2m).
\]
It follows that the last bound is at most
\[
 (66+8\cdot400^2)\log(2m)
 \le c_{\mathrm{cl}}\log(2m).
\]
Only the number of partitions enters this estimate.  The number of
coordinates within a batch, and the eventual modulus, may depend on
the entire roster.

Finally put $n=\prod_{c\in\mathcal C}N_c$.  The residue map
\[
 \mathbb Z/n\mathbb Z\longrightarrow G,
 \qquad x\longmapsto(x\bmod N_c)_{c\in\mathcal C},
\]
is an isomorphism of additive groups.  For an explicit proof of
surjectivity, $n/N_c$ is nonzero modulo the prime $N_c$, so some
multiple $e_c$ of it is $1$ modulo $N_c$; the same $e_c$ is zero in
all other coordinates.  A prescribed tuple $(x_c)$ is therefore the
image of $\sum_c x_ce_c$.  The kernel consists of integers divisible
by every $N_c$, hence by their product, proving injectivity.  Transport
$I_\alpha,J_\alpha,W$, and all $v_p$ through this isomorphism.  All
sum, difference, and translation statements are preserved, and every
prime factor of $n$ satisfies \textup{(i)}.

For $m=1$ the first-coordinate construction still uses the positive
number $a$ of coordinates, while no Type batches are needed.  The
separation constraints are vacuous, and the diagonal is available
whenever all first totals are nonzero.  Thus this case uses the same
proof, without an empty product of moduli.  Empty and singleton
rosters have no ordered distinct pairs and were also included above.
This completes the proof.
\end{proof}

\subsection{Interpretation at a cut}

Suppose a prefix test selects $\alpha$ and knows an integer $i$, while
a suffix test independently selects $\beta$ and knows an integer $j$.
Their clock requirements are
\[
 i\bmod n\in I_\alpha,\qquad j\bmod n\in J_\beta.
\]
If $i+j=\ell$ is fixed across all successful cuts under consideration,
Lemma~\ref{lem:clock} applies to every such pair of tests.  In the
diagonal case their tags must agree.  In the off-diagonal case every
successful pair uses the same prefix tag $\alpha$ and the same suffix
tag $\beta$.
In the episode applications below,
\[
 i=k-t,\qquad j=q-k,\qquad \ell=q-t.
\]
The common total is available only after an independent end guard
has established the actual anchor $q$.  The clock does not itself
establish the anchor or the validity of other metadata.

The tags will include a claimed input $X$ and output $Y$ in a vector
space, with the prefix checking $x=X(\alpha)$ and the suffix checking
$y=Y(\beta)$.  If the unique edge is off-diagonal, every successful
transfer therefore has the same input and output, and is consistent
with the affine map
\[
 v\longmapsto vT+Y(\beta)-X(\alpha)T
\]
for any specified linear part $T$.  This explains why the exclusive
alternative in Lemma~\ref{lem:clock} is sufficient for the affine
updates used below.  An empty graph simply permits no successful
transfer.

For availability, if $\ell\notin W$, every desired diagonal tag
$\alpha$ has some decomposition $\ell=i+j$ with
$i\in I_\alpha$ and $j\in J_\alpha$.  Thus a family of candidate cuts
whose first residues exhaust $\mathbb Z/n\mathbb Z$ contains a cut
meeting both requirements for that tag, provided its other checks
hold throughout the family.  Consecutive $n$ cuts have this property,
as does a progression of $n$ cuts whose step is coprime to $n$.
Indeed multiplication by a coprime step permutes the residue classes.

Lastly, suppose integer offsets $p$ have been assigned the roster
residues $v_p$, and $q,t,\gamma$ are fixed integers.  Then
\[
 q-(t+\gamma+p-p')\in W
 \quad\Longleftrightarrow\quad
 (q-t-\gamma)+v_{p'}-v_p\in W,
\]
where membership uses residues modulo $n$.  Equation~\eqref{eq:clock-roster}
bounds the number of ordered pairs satisfying this condition.
For a fixed finite set of translates $\Gamma$, the union over
$\gamma\in\Gamma$ has at most
$|\Gamma|c_{\mathrm{cl}}\log(2m)$ such pairs.  Both $q$ and each
$\gamma$ are held fixed while applying the roster estimate; choosing
a usable offset afterwards does not alter that bound.
